(1) 在两平板 \(y=\pm h\) 间的定常层流速度为 \(u(y)=U_{\max}(1-y^2/h^2)\)。被动溶质浓度满足 \(\partial_t c+u(y)\partial_xc=D(\partial_{xx}+\partial_{yy})c\)，\(D>0\)，两壁不透质 \(\partial_yc=0\)。在横向扩散时间 \(t\gg h^2/D\)、纵向慢变条件下，推导平均速度 \(\bar U\)，并通过渐近胞元修正 \(c=\bar c+\chi(y)\partial_x\bar c+\cdots\) 得到决定修正函数的Neumann方程及零均值条件。
(2) 严格求 \(\chi'(y)\)，利用胞元解推导精确领先阶有效纵向扩散系数 \(D_{\rm eff}\)，要求实际计算截面多项式积分、保留准确系数。
(3) 独立使用纵向随体坐标 \(X=x-\bar Ut\) 中的溶质质量一阶与二阶矩，说明长期平均剖面所遵从的方程与初始有限质量脉冲的方差增长率，解释 \(D_{\rm eff}\) 必大于分子扩散系数 \(D\)。
(4) 用 \(Pe=\bar Uh/D\) 表示增强弥散相对大小，检验 \(Pe\to0\) 与 \(h\to0\) 的极限、有效系数的量纲；说明对横向尚未混合的早期溶质团为何不能应用上述常系数Fick扩散。
